\documentclass[12pt]{article}
\usepackage[T1]{fontenc}
\usepackage[utf8]{inputenc}
\usepackage{lmodern}
\usepackage{textcomp}
\usepackage{enumitem}
\usepackage{geometry}
\usepackage{graphicx}
\usepackage{amsmath, amssymb}
\geometry{margin=1in}
\begin{document}
\begin{center}
Economics - Graduate Level Assessment 25 \
Total Marks: 40 \
Answer all questions.
\end{center}

\section*{Multiple Choice (10 marks)}
\begin{enumerate}[label=\arabic*.]
\item When a firm is earning a normal profit from the production of a good, it is true that
\begin{enumerate}[label=(\alph*)]
    \item total revenues from production are equal to implicit costs.
    \item total revenues from production are equal to explicit costs.
    \item explicit costs are higher than implicit costs.
    \item total revenues from production are equal to the sum of explicit and implicit costs.
    \item total revenues from production are lower than the sum of explicit and implicit costs.
\end{enumerate}
\item Which of the following will shift the supply curve for textbooks to the left?
\begin{enumerate}[label=(\alph*)]
    \item An increase in printing costs
    \item Expectations of future surpluses
    \item A decrease in the demand for a substitute in production
    \item An increase in the literacy rate
    \item A decrease in the cost of labor for bookbinding
\end{enumerate}
\item Which of the following statements is correct in regard to the federal budget deficit and the federal debt?
\begin{enumerate}[label=(\alph*)]
    \item The debt is the sum of all past and future deficits.
    \item When the debt is positive the deficit decreases.
    \item When the deficit is negative the debt decreases.
    \item The deficit is the accumulation of past debts.
    \item When the debt is negative the deficit decreases.
\end{enumerate}
\item The profit-maximizing equilibrium for a monopolistically competitive firm leaves marginal cost below price. Explain why this is inefficient from a societal per-spective .
\begin{enumerate}[label=(\alph*)]
    \item Society does not get enough output.
    \item There is perfect competition in the market, leading to an efficient outcome.
    \item Monopolistically competitive firms produce at the minimum of the average cost curve.
    \item There is a socially efficient allocation of resources.
    \item The firm's profits are minimized, leading to an equitable distribution of resources.
\end{enumerate}
\item Which of the following is true?
\begin{enumerate}[label=(\alph*)]
    \item TFC = AVC at all levels of output.
    \item AVC + AFC = TC.
    \item TFC = TC at all levels of output.
    \item MC = AVC + AFC.
    \item AFC = TC - AVC.
\end{enumerate}
\end{enumerate}

\section*{True / False (10 marks)}
\textit{Answer True or False}
\begin{enumerate}[label=\arabic*.]
\setcounter{enumi}{5}
\item True or False: A progressive income tax adds an element of built-in stability to an economy by adjusting tax rates based on income levels.
\item True or False: Contractionary monetary policy tends to increase interest rates, leading to falling Treasury bond prices.
\item True or False: An increase of $100 in autonomous private investment results in a $1,000 shift to the right in aggregate demand with a marginal propensity to consume of 0.90.
\item True or False: Selling $29 million in government securities will cause the money supply to increase by $290 million through the money multiplier effect.
\item True or False: The value of the metal in each United States coin is significantly lower than its designated monetary value.
\end{enumerate}

\section*{Long Form Response (20 marks)}
\textit{Answer each question with a clear and complete explanation.}
\begin{enumerate}[label=\arabic*.]
\setcounter{enumi}{10}
\item Discuss the fundamental differences between the short-run and long-run Phillips curves, focusing on their shapes, implications for inflation and unemployment, and economic policy considerations.
\item Discuss how changes in the required reserve ratio affect the money supply in an economy, including the mechanisms through which this impact occurs.
\end{enumerate}

\vfill
\noindent\textit{End of Paper}
\end{document}